\part{Comparison theorems and degenerations}
\label{part:comparisons}

The final part assembles the promised dictionaries.  Vertical divisors,
model metrics, CLD currents, Fubini--Study functions, filtrations, test
curves, geodesic rays, and complex singularity data are compared on their
common domains.  Separate chapters treat Monge--Amp\`ere measures, Legendre
transforms, hybrid convergence, functoriality, and stability.  The research
chapters retain only conclusions with a recorded proof status; unresolved
attempts remain problems with their exact obstruction rather than being
presented as theorems.

The part also changes scale twice.  The normalized valuation link of a
pointed smooth germ leads to local nef b-divisors, ideal functions, and local
intersection measures.  The compact K\"ahler discussion then returns to
pseudoeffective classes and asks how the Darvas--Xia--Zhang inverse system is
seen on the Mesquita--Piccione tropical spectrum.  In both directions, the
point is to identify proved interfaces and missing arrows, not to force all
psh notions into a single definition.

\paragraph{Prerequisites.}
This is not intended as an independent first course.  For BFJ--CLD
comparisons the reader should know
\Cref{chap:metrics,chap:cld-currents,chap:bfj-model-forms,chap:bfj-ma}; for
BJ, filtrations, and rays, the minimum is
\Cref{chap:bj-fs-pl-psh,chap:bj-energy,chap:rays-test-curves,chap:dxz-bj-comparison}.
The local chapter uses valuation and b-divisor material from
\Cref{chap:valuations,app:intersection-bdivisors}; the pseudoeffective chapter
uses both clusters of \Cref{part:complex}.  Readers entering at a comparison
chapter should also keep \Cref{app:normalizations,app:provenance} nearby.

\paragraph{Conceptual clusters.}
The dictionary cluster begins with \Cref{chap:master-dictionary}, which
distinguishes equivalences, one-way constructions, and analogies by attaching
a domain signature to every arrow.  Monge--Amp\`ere measures are compared in
\Cref{chap:compare-ma}, and filtrations, FS metrics, test curves,
Duistermaat--Heckman measures, and rays in \Cref{chap:rays-filtrations}.  The
degeneration cluster consists of
\Cref{chap:hybrid-spaces,chap:functoriality-trace,chap:ytd-bridge}; hybrid
limits, pullback and trace, and the conditional stability bridge are kept in
separate chapters so that a convergence theorem cannot silently supply a
missing entropy or coercivity statement.

The frontier cluster is
\Cref{chap:verified-extensions,chap:local-valuative-psh,chap:pseff-mp-dxz,chap:major-open-problems}.
The first chapter records audited extensions, reductions, and
counterexamples.  The local chapter develops the established valuative core
and itemizes the additional data required for a local BJ-type theory.  The
pseudoeffective chapter compares the MP and DXZ constructions on the domains
for which the ideal-valuative transform and its operations have been
verified, while keeping the big non-nef and zero-volume boundaries visible.
The final chapter collects the remaining source-dated problems and records a
source-scoped verification status for each; a queued item is not counted as
an attempted solution.

\paragraph{Reading routes.}
The \emph{short route} to the established global dictionary is
\Cref{chap:master-dictionary,chap:compare-ma,chap:rays-filtrations}.  The
\emph{standard route} adds
\Cref{chap:hybrid-spaces,chap:functoriality-trace,chap:ytd-bridge} and then
reads the status ledger in \Cref{chap:verified-extensions}.  The
\emph{reference route} reads the entire part together with
\Cref{app:normalizations,app:provenance}; it is the only route that includes
the local theory, the MP--DXZ pseudoeffective boundary, and the complete
open-problem ledger.  A reader interested only in local positivity may pass
from \Cref{chap:cld-examples} to
\Cref{chap:verified-extensions,chap:local-valuative-psh}; a reader interested
in complex degenerations may pass from \Cref{part:complex} to
\Cref{chap:rays-filtrations,chap:hybrid-spaces,chap:ytd-bridge,chap:pseff-mp-dxz}.

\paragraph{Main outcomes and the six-way interface.}
The part shows how to ask a comparison question correctly.  BFJ contributes
model-skeletal psh functions and their global nontrivially valued
Monge--Amp\`ere theory.  BJ contributes homogeneous FS functions,
filtrations, b-divisors, and the trivially valued energy theory.  CLD
contributes local forms and currents; the Gubler--Rabinoff approach adds
harmonic charts and curve tests.  The complex theory contributes rays, test
curves, multiplier ideals, and radial slopes, while the MP and DXZ
constructions provide two transcendental inverse-system realizations.  Their
intersection is largest on model, PL, divisorial, projective, or
finite-energy loci, depending on the arrow.  Outside those loci,
\Cref{chap:master-dictionary,chap:verified-extensions,chap:pseff-mp-dxz}
records whether the book has an equivalence, a one-way construction, a
conditional reduction, a counterexample, or only an analogy.

The rational toric numerical-dimension comparison is deliberately checked at
both interfaces: the full Cao--Guenancia--character calculation in
\Cref{chap:transcendental-okounkov-dh} and the separately frozen statement and
proof in \Cref{chap:verified-extensions} have independent exact-current
production receipts.  This is a source synthesis, not a novelty claim, and it
does not include irrational real toric classes.

\paragraph{Normalization and status warnings.}
Model intersection measures are compared before probability normalization;
division by volume occurs only when that volume is positive.  Hybrid slopes
must respect the squared-norm convention for complex potentials, and
divisorial formulas must distinguish integral orders from valuations
normalized by a degeneration parameter.  Restriction is not trace, weak
convergence is not strong convergence, and agreement on divisorial points is
not automatically equality on an entire nonmetrizable spectrum.  Finally,
the words \emph{verified}, \emph{conditional}, \emph{disproved}, and
\emph{open} are used as mathematical status labels; their provenance is
audited in \Cref{app:provenance}.
